%
% Spezielle komplexe Funktionen
%
\subsubsection{Spezielle komplexe Funktionen}

%
% produkt: Produktdarstellungen
%
Polynome sind bis auf einen Vorfaktor durch die Nullstellen
festgelegt.
Ob dies auch für beliebige ganze Funktionen gilt, diskutieren
\emph{Simon Köpfli}
\index{Kopfli, Simon@Köpfli, Simon}%
\index{Simon Kopfli@Simon Köpfli}%
und
\emph{Florian von Wyl}
\index{von Wyl, Florian}%
\index{Florian von Wyl}%
in Kapitel~\ref{chapter:produkt}.
Da es ganze Funktionen mit unendlich vielen Nullstellen gibt, muss
man dazu unendlichen Produkten einen Sinn geben.
\index{Produkt, unendlich}%
Dabei treten zunächst Konvergenzfragen als Hindernis auf, die durch
die weierstraßschen Elementarfaktoren beantwortet werden.
Dies genügt aber nicht.
Die Exponentialfunktion $e^{g(z)}$ einer ganzen Funktion $g(z)$ hat
keine Nullstellen, die möglicherweise verschiedenen Funktionen $f(z)$
und $e^{g(z)}f(z)$ haben daher die gleichen Nullstellen, es können
aber im Allgemeinen nicht beide Funktionen Polynome sein.

%
% basel: Basel-Problem
%
Das Basel-Problem hat die besten Mathematiker des 17.~und 18.~Jahrhunderts
herausgefordert.
\index{Basel-Problem}%
Erst Leonhard Euler fand die korrekte Summe der reziproken Quadrate der
\index{Euler, Leonhard}%
natürlichen Zahlen.
\emph{Jan Gachnang}
und
\emph{Etienne Schafflützel}
präsentieren in Kapitel~\ref{chapter:basel} Eulers klassische Lösung,
die allerdings modernen Standards nicht genügt.
Strenge Beweise werden ebenfalls dargelegt.

%
% gamma: Gamma-Funktion
%
Die Fakultät ist die prototypische rekursiv definierte Funktion, sie
tritt in der Taylor-Reihe und in der Kombinatorik immer wieder auf,
und sie wird auch in den Kapiteln~\ref{chapter:zeta} und \ref{chapter:hankel}
verwendet.
\emph{Lukas Buchli}
\index{Lukas Buchli}%
\index{Buchli, Lukas}%
und
\emph{Roman Weber}
\index{Roman Weber}%
\index{Weber, Roman}%
zeigen, wie die gleiche rekursive Struktur auch in den Volumina 
geradedimensionaler Kugeln auftritt.
Sie leiten die Formel
\[
V_n(r)
=
\frac{\pi^{n/2}}{(n/2)!}r^n
\]
für das Volumen der $n$-dimensionalen Kugel her.
Die Formel funktioniert natürlich nur für gerade $n$, doch wenn man
die Fakultät auch für beliebige reelle Zwischenwerte definieren könnte,
würde die Formel auch für ungeraddimensionale Kugeln funktionieren.
Wie der Satz von Bohr-Mollerup diese Erweiterung zur Gamma-Funktion
\index{Gamma-Funktion}%
möglich macht, zeigen sie in Kapitel~\ref{chapter:gamma}.

%
% bessel: Bessel-Funktionen
%
TODO: bessel

%
% hankel: Hankel
%
TODO: hankel

%
% zeta: Die Summe der natürlichen Zahlen
%
Die Berechnung von Werten der Zeta-Funktion ist nicht immer einfach.
Euler berechnet für das Basel-Problem erfolgreich den Wert
$\zeta(2)=\frac{\pi^2}{6}$.
Kapitel~\ref{chapter:zeta} von
\emph{Melissa Haumüller}
\index{Melissa Haumüller}%
\index{Haumüller, Melissa}%
zeigt, wie sich mit der analytischen Erweiterung der $\zeta$-Funktion
der auf den ersten Blick unsinnige Wert
$\zeta(-1) = 1 + 2 + 3 +\dots = -\frac{1}{12}$ recthfertigen lässt.
